\documentclass[compress, 8pt, aspectratio=169]{beamer}

\usepackage{presentationtemplate}
\usepackage[askip=3mm, bskip=3mm]{terminal}
\usepackage[linenosfontsize=\tiny, askip=3mm, bskip=3mm]{mylisting}
\tcbuselibrary{hooks}
\usepackage{tcolorbox}
\usepackage{csquotes}
\usepackage{tabularray}

\newtcolorbox{task}{
    colback=yellow!50!white,
    boxrule=0.02cm,
    colframe=black,
    sharp corners,
    left=0mm,
    right=0mm,
    top=0mm,
    bottom=0mm,
    before upper={\textbf{Задание}:\:},
}

\title{Стандартная библиотека: алгоритмы}

\begin{document}

  \frame[plain]{\titlepage}

  \begin{frame}[fragile]

    \frametitle{Алгоритмы из STL}

    \begin{itemize}

      \item \textbf{Алгоритм}: шаблонная функция
        из заголовочного файла \texttt{<algorithm>}.
      \item Не хранит данные: итераторы на данные
        передаются аргументами.
      \item Работает с любым контейнером:
        \texttt{array}, \texttt{vector}, \texttt{list}, обычный массив.
      \item Условие и преобразование передаются отдельным аргументом,
        функцией (например, адресом именованной функции, замыканием
        лямбда-функции и т.д.).

    \end{itemize}

  \end{frame}

  \begin{frame}[fragile]

    \frametitle{Типичная сигнатура алгоритма}

    Большинство алгоритмов принимают диапазон итераторов
    \texttt{[first, last)} и возвращают итератор или значение:

    \begin{myinplacelisting}[minted language=cpp]
template <class InIt, class OutIt>
OutIt copy(InIt first, InIt last, OutIt d_first);
    \end{myinplacelisting}

    \begin{itemize}

      \item \texttt{first}, \texttt{last}: обрабатываемый диапазон;
      \item \texttt{d\_first}: начало диапазона, куда пишется результат;
      \item результат \textemdash \space итератор на элемент, следующий за
        последним обработанным.

    \end{itemize}

  \end{frame}

  \begin{frame}[fragile]

    \frametitle{Категории \texttt{<algorithm>}}

    \begin{tblr}{
      colspec = {|X[0.30,l]|X[0.26,l]|X[0.44,l]|},
      row{1} = {font=\bfseries},
      hlines, vlines,
      rowsep = 0.8ex,
      colsep = 1.0ex,
    }
      Категория (cppreference) & Алгоритм & Что делает \\
      Copy operations & \texttt{std::copy} & Копирует диапазон в другое место. \\
      Transformation & \texttt{std::transform} & Применяет функцию к каждому элементу. \\
      Generation & \texttt{std::fill} & Заполняет диапазон значением. \\
      Removing & \texttt{std::remove} & \enquote{Логически убирает} значения из диапазона. \\
      Order-changing & \texttt{std::reverse} & Переворачивает порядок элементов. \\
      Search operations & \texttt{std::find} & Находит первый элемент, равный значению. \\
      Sorting operations & \texttt{std::sort} & Сортирует диапазон. \\
      Minimum/maximum operations & \texttt{std::max\_element} & Находит наибольший элемент, возвращает итератор. \\
    \end{tblr}

  \end{frame}

  \begin{frame}[fragile]

    \frametitle{\texttt{std::copy}}

    \myinputlisting[minted language=cpp]{Presentations/19-algorithms/}{copy.cpp}

    \begin{terminalwindow}
!\shellcommand{c++ -std=c++23 copy.cpp -o copy}!
!\shellcommand{./copy}!
1 2 3
    \end{terminalwindow}

  \end{frame}

  \begin{frame}[fragile]

    \frametitle{\texttt{std::copy}: что возвращается}

    Возвращает итератор на элемент, следующий за последним
    скопированным, его можно передать дальше:

    \begin{myinplacelisting}[minted language=cpp]
auto last = std::copy(src.begin(), src.end(),
                      dst.begin());
// last == dst.end()
    \end{myinplacelisting}

    \begin{task}
      Что вернёт \texttt{std::copy} для пустого диапазона?
      Как можно воспользоваться возвращённым итератором?
    \end{task}

  \end{frame}

  \begin{frame}[fragile]

    \frametitle{\texttt{std::transform}}

    \myinputlisting[minted language=cpp]{Presentations/19-algorithms/}{transform.cpp}

  \end{frame}

  \begin{frame}[fragile]

    \frametitle{\texttt{std::transform}: вывод программы}

    Вывод программы с листинга на прошлом слайде:

    \begin{terminalwindow}
!\shellcommand{c++ -std=c++23 transform.cpp -o transform}!
!\shellcommand{./transform}!
1 4 9 16
    \end{terminalwindow}

  \end{frame}

  \begin{frame}[fragile]

    \frametitle{\texttt{std::fill}}

    Присваивает одно и то же значение каждому элементу диапазона.

    \myinputlisting[minted language=cpp]{Presentations/19-algorithms/}{fill.cpp}

    \begin{terminalwindow}
!\shellcommand{c++ -std=c++23 fill.cpp -o fill}!
!\shellcommand{./fill}!
7 7 7 7 7
    \end{terminalwindow}

  \end{frame}

  \begin{frame}[fragile]

    \frametitle{\texttt{std::remove}}

    \enquote{Логически убирает} все вхождения значения:
    оставшиеся элементы сдвигаются вперёд, возвращается итератор
    на новую границу.
    Размер контейнера не меняется.

    \myinputlisting[minted language=cpp]{Presentations/19-algorithms/}{remove.cpp}

  \end{frame}

  \begin{frame}[fragile]

    \frametitle{\texttt{std::remove}: вывод программы}

    Вывод программы с листинга на прошлом слайде:

    \begin{terminalwindow}
!\shellcommand{c++ -std=c++23 remove.cpp -o remove}!
!\shellcommand{./remove}!
1 3
size: 5
    \end{terminalwindow}

  \end{frame}

  \begin{frame}[fragile]

    \frametitle{Идиома erase-remove}

    Чтобы элемент реально исчез из контейнера,
    результат \texttt{std::remove} передаётся
    в \texttt{container::erase}.

    \myinputlisting[minted language=cpp]{Presentations/19-algorithms/}{erase-remove.cpp}

  \end{frame}

  \begin{frame}[fragile]

    \frametitle{Идиома erase-remove: вывод программы}

    Вывод программы с листинга на прошлом слайде:

    \begin{terminalwindow}
!\shellcommand{c++ -std=c++23 erase-remove.cpp -o erase-remove}!
!\shellcommand{./erase-remove}!
1 3
size: 2
    \end{terminalwindow}

    \begin{task}
      Подумайте, зачем потребовался такой подход к удалению элементов?
      В чем его преимущество между обычным удалением в цикле?
    \end{task}

  \end{frame}

  \begin{frame}[fragile]

    \frametitle{\texttt{std::reverse}}

    Меняет порядок элементов в диапазоне на обратный.

    \myinputlisting[minted language=cpp]{Presentations/19-algorithms/}{reverse.cpp}

    \begin{terminalwindow}
!\shellcommand{c++ -std=c++23 reverse.cpp -o reverse}!
!\shellcommand{./reverse}!
4 3 2 1
    \end{terminalwindow}

  \end{frame}

  \begin{frame}[fragile]

    \frametitle{\texttt{std::find}}

    Находит первый элемент, равный значению,
    и возвращает итератор на него.
    Возвращает \texttt{end}, если не нашёл.

    \myinputlisting[minted language=cpp]{Presentations/19-algorithms/}{find.cpp}

    \begin{task}
      Что вернёт \texttt{std::find}, если равного элемента нет?
    \end{task}

  \end{frame}

  \begin{frame}[fragile]

    \frametitle{\texttt{std::find}: вывод программы}

    Вывод программы с листинга на прошлом слайде:

    \begin{terminalwindow}
!\shellcommand{c++ -std=c++23 find.cpp -o find}!
!\shellcommand{./find}!
found: 7
    \end{terminalwindow}

  \end{frame}

  \begin{frame}[fragile]

    \frametitle{\texttt{std::sort}}

    Сортирует элементы диапазона, после вызова
    диапазон отсортирован.

    \myinputlisting[minted language=cpp]{Presentations/19-algorithms/}{sort.cpp}

    \begin{terminalwindow}
!\shellcommand{c++ -std=c++23 sort.cpp -o sort}!
!\shellcommand{./sort}!
1 2 3 4 5
    \end{terminalwindow}

  \end{frame}

  \begin{frame}[fragile]

    \frametitle{\texttt{std::max\_element}}

    Находит наибольший элемент диапазона и возвращает
    итератор на него, так же как \texttt{std::find}.

    \myinputlisting[minted language=cpp]{Presentations/19-algorithms/}{max-element.cpp}

    \begin{terminalwindow}
!\shellcommand{c++ -std=c++23 max-element.cpp -o max-element}!
!\shellcommand{./max-element}!
max: 5
    \end{terminalwindow}

  \end{frame}

  \begin{frame}[fragile]

    \frametitle{Собственный алгоритм: find\_even}

    Стандартный алгоритм \textemdash \space это шаблон,
    который \enquote{проходит} диапазон по итераторам
    и возвращает итератор.
    Такой алгоритм можно разработать самостоятельно.

    \myinputlisting[minted language=cpp, minted options app={fontsize=\footnotesize}]{Presentations/19-algorithms/}{find-even.cpp}

  \end{frame}

  \begin{frame}[fragile]

    \frametitle{find\_even: вывод программы}

    Вывод программы с листинга на прошлом слайде:

    \begin{terminalwindow}
!\shellcommand{c++ -std=c++23 find-even.cpp -o find-even}!
!\shellcommand{./find-even}!
first even: 4
    \end{terminalwindow}

  \end{frame}

  \begin{frame}[fragile]

    \frametitle{Тот же результат даёт \texttt{std::find\_if}}

    \myinputlisting[minted language=cpp]{Presentations/19-algorithms/}{find-if.cpp}

  \end{frame}

  \begin{frame}[fragile]

    \frametitle{\texttt{std::find\_if}: вывод программы}

    Вывод программы с листинга на прошлом слайде:

    \begin{terminalwindow}
!\shellcommand{c++ -std=c++23 find-if.cpp -o find-if}!
!\shellcommand{./find-if}!
first even: 4
    \end{terminalwindow}

  \end{frame}

  \begin{frame}[fragile]

    \frametitle{Композиция стандартных алгоритмов}

    Функции собирают из стандартных алгоритмов:
    \texttt{erase\_if} = \texttt{std::remove\_if} + \texttt{c::erase}.
    Так же устроен \texttt{std::erase\_if} из C++20.

    \myinputlisting[minted language=cpp, minted options app={fontsize=\footnotesize}]{Presentations/19-algorithms/}{erase-if.cpp}

  \end{frame}

  \begin{frame}[fragile]

    \frametitle{\texttt{erase\_if}: вывод программы}

    Вывод программы с листинга на прошлом слайде:

    \begin{terminalwindow}
!\shellcommand{c++ -std=c++23 erase-if.cpp -o erase-if}!
!\shellcommand{./erase-if}!
1 3 5
    \end{terminalwindow}

  \end{frame}

\end{document}
